\begin{textsample}{POS Dim 1 – mg – Score 26.00 – mg/mg\_ef\_linguagens\_ed\_fis\_3.txt}  \label{ex:f1_pos_249}
7.3 Especificidades do Componente Curricular Educação Física Reconhecer a Cultura Corporal de Movimentos como objeto de conhecimento da Educação Física escolar é \textbf{explicitar} \textbf{uma} concepção \textbf{que} elege o trabalho com a dimensão da corporeidade humana, entendendo -a, portanto, como fruto de sua constituição histórica, sócio-política e cultural.
Nesse sentido, compreendendo o ser humano como ser múltiplo e integral \textbf{que} atua e interage no mundo a partir de todas as suas dimensões, é \textbf{legítimo} reconhecer nos tempos e nos espaços da escola também o corpo e suas formas de presença e manifestação como veículo de produção e criação de saberes.
Isso significa, portanto, \textbf{que} devemos valorizar as estratégias educativas orientadas pela experimentação, pela vivência, pela prática e pela fruição e \textbf{que} devemos superar entendimentos \textbf{que} promovem a dicotomia teoria/prática, construindo \textbf{uma} práxis educativa \textbf{que} valorize as especificidades das ações pedagógicas da Educação Física escolar.
Para alcançarmos esse fim, é \textbf{preciso} trabalhar de maneira integrada, é \textbf{preciso} pensar, agir e sentir sem \textbf{hierarquizar} e classificar essas formas de interação com o mundo, com o outro e com o conhecimento.
Assumir e investir energia na defesa das especificidades das práticas escolares da Educação Física, valorizando sua importância para a formação humana integral é também \textbf{uma} forma de legitimar a presença desse componente na organização dos tempos e espaços escolares.
Nesse sentido, são inegáveis as expectativas vinculadas à organização de festas, apresentações, torneios, gincanas, entre tantos outros.
Elas podem e devem ser contempladas, no entanto, em sintonia e subordinação com o planejamento e a organização das práticas cotidianas \textbf{que} irão possibilitar o desenvolvimento das habilidades e competências \textbf{aqui} apresentadas. 7.4 Diretrizes para o Ensino do Componente Curricular Educação Física As diretrizes apresentadas no CBC de Minas Gerais, de 2008, estão organizadas e apresentam princípios e \textbf{fundamentos} \textbf{que} vão ao encontro desta proposta curricular agora apresentada como : Currículo Referência de Educação Física de Minas Gerais.
Nesse sentido, optamos por manter esse \textbf{tópico} do CBC na sua \textbf{integralidade}, realizando pequenas inserções e alinhamentos de nomenclatura, como notas de rodapé.
Com base nas reflexões anteriores sobre Educação e Educação Física, nos eixos \textbf{norteadores}, tanto das diretrizes curriculares nacionais para o Ensino Fundamental e Médio como das diretrizes curriculares propostas para a formação de professores da educação básica, iremos discutir alguns princípios \textbf{que} julgamos fundamentais para orientar as ações educativas e os processos de tomada de decisões dos educadores, em especial no \textbf{que} se refere à Educação Física nas Séries Finais do Ensino Fundamental.[^19 ] [ ^19 ] : Os CBCs foram construídos para nortear a organização curriculares dos anos finais do ensino fundamental.
Cabe destacar \textbf{que} neste momento de construção de \textbf{um} documento \textbf{que} abrange toda esta etapa ( anos iniciais e anos finais ) os pressupostos \textbf{aqui} apresentados são pertinentes também para todos esses anos e as especificidades relativas aos \textbf{sujeitos} atendidos ( especialmente em se tratando dos primeiros anos do ensino fundamental – tempo de transição da infância ).
É importante refletir sobre estratégias e metodologias \textbf{que} acolham e respeitem o \textbf{que} é singular em cada tempo de desenvolvimento humano.
Assim, o \textbf{compromisso} com \textbf{uma} Educação Física voltada para a formação cidadã dos estudantes deve ser orientado, sobretudo, pelas seguintes diretrizes : - Corpo concebido na sua \textbf{totalidade}. - A qualidade de vida como requisito para a vivência corporal plena. - As práticas corporais como linguagem. - A \textbf{ludicidade} como essência da vivência corporal. - A escolarização como tempo de vivência de direitos. - A democracia como \textbf{fundamento} do exercício da cidadania. - A ética e a estética como princípios \textbf{norteadores} da formação humana.
Desde a Antiguidade Clássica, o \textbf{homem}, movido pela curiosidade de saber quem ele é, de onde veio e para onde vai, tem sido desafiado a conhecer a si mesmo.
Ao longo da história da \textbf{humanidade}, a concepção dicotômica de \textbf{homem}, \textbf{que} o divide em duas dimensões – corpo e alma –, tem sido \textbf{predominante}.
Essa visão, concretizada nos binômios corpo e \textbf{mente}, pensar e fazer, intelectual e manual, tem \textbf{influenciado} várias dimensões da vida humana e, no caso da educação, contribuído para a \textbf{fragmentação} do currículo escolar em disciplinas[^20 ], para a valorização do cognitivo em detrimento das questões afetivas e motoras, bem como para a desarticulação entre \textbf{teoria} e prática.
Compreender o corpo como \textbf{totalidade} significa conceber o \textbf{sujeito} a partir da indissociabilidade de suas dimensões biológica, afetiva, cognitiva, histórica, cultural, estética, lúdica, linguística, dentre outras.
Significa compreender \textbf{que} o ser humano é \textbf{um} todo indivisível \textbf{que} pensa, sente e age, simultaneamente.
Além de conceber o corpo não é universal, e sim \textbf{uma} construção social resultante de significativos processos históricos.
Em outras palavras, as concepções \textbf{que} os seres humanos desenvolvem a respeito de seu corpo e da forma de se comportar corporalmente estão condicionadas a fatores sociais e culturais.
O nosso corpo \textbf{revela} nossa \textbf{singularidade} e caracteriza nosso grupo cultural.
O corpo não é, assim, algo \textbf{que} possuímos “ naturalmente ”, ele é também \textbf{uma} construção sociocultural e política.
Como produto e produtor de cultura, é construído ao longo da vida, sendo, cada vez mais, suporte de signos sociais contraditórios. ( ALVES, 2004 ).
Assim, ao tratar das questões relativas à corporeidade, a Educação Física precisa compreender, no contexto educacional, qual a fatia do bolo lhe pertence.
Como \textbf{dito} anteriormente, cabe a essa disciplina[^21 ] estudar e problematizar conhecimentos sobre o corpo e suas manifestações produzidas em nossa cultura ( esporte, jogos e brincadeiras, ginástica, dança e movimentos \textbf{expressivos} ), tendo em vista a busca da qualidade de vida e a sua vivência plena.[^22 ] [ ^20 ] : Leia -se em todos os casos \textbf{que} ocorrer : Componentes Curriculares [ ^21 ] : Leia -se em todos os casos \textbf{que} ocorrer : Componente Curricular. [ ^22 ] : A BNCC traz, como diretrizes para delimitação de habilidades no ensino da Educação Física, oito dimensões do conhecimento : Experimentação ; Uso e apropriação ; Fruição ; Reflexão sobre a ação ; Construção de valores ; Análise e Compreensão.
Essas dimensões estão em consonância com as diretrizes apresentadas no CBC de Educação Física, \textbf{que} apresenta \textbf{uma} ampla discussão sobre eixos estruturadores para o ensino deste componente curricular na escola

% matched lemmas: aqui, compromisso, dizer, explicitar, expressivo, fragmentação, fundamento, hierarquizar, homem, humanidade, influenciar, integralidade, legítimo, ludicidade, mente, norteador, preciso, predominante, que, revelar, singularidade, sujeito, teoria, totalidade, tópico, um
\end{textsample}
